\begin{foldable}
    Our study explores \textit{black-box data-free knowledge distillation}, a setting where a student model learns from a teacher's outputs without access to its internal parameters or original training data.
    While not apparent, we would like to raise certain ethical and societal concerns related to model extraction, information leakage, and bias propagation.
\end{foldable}

\begin{foldable}
    Because the student learns entirely from the teacher's responses, repeated querying could approximate the teacher's decision boundaries, potentially enabling unauthorized reconstruction of proprietary or commercial models \citep{tramer2016stealing,orekondy2019knockoff,jagielski2020high}.
    Large-scale replication of closed models could discourage openness in AI research and reduce incentives for model providers to offer public APIs.
    This problem could be mitigated by limiting the queries to the teacher, and through model watermarking or output perturbation that make functional copying more detectable or less precise \citep{jia2021entangled,wang2024defense}.
\end{foldable}

\begin{foldable}
    Besides, even though no real data are used, teacher models may implicitly memorize sensitive information that can be exposed through their responses \citep{carlini2021extracting}.
    In extreme cases, this could lead to indirect leakage of personal or confidential data.
    Such leakage risks eroding public trust in machine learning systems and may breach data protection norms if human-derived models are involved.
    As a recommendation, using privacy-preserving training techniques, such as differential privacy \citep{flemings2024differentially}, and conducting ethical audits of teacher outputs can reduce the likelihood of revealing sensitive patterns.
\end{foldable}

\begin{foldable}
    In data-free KD, the student may be trained using only the teacher's labels, thus it can directly inherit and sometimes amplify distillation-induced biases in the teacher's decision boundaries \citep{hooker2020characterising,ahn2022knowledge}.
    Those amplified disparities can disproportionately harm marginalized groups when distilled models are deployed in downstream settings \citep{buolamwini2018gender}.
    Mitigation strategies include using fairness-aware distillation and audit methods \citep{chai2022fairness,masroor2024fair}.
\end{foldable}

\begin{foldable}
    Our aim is to enable ethical and resource-efficient research on knowledge distillation, not to facilitate extraction or imitation of closed systems.
    We therefore restrict our experiments to openly licensed teacher models and open-source our works to encourage transparent, auditable research practices.
\end{foldable}
